% !TEX root=../main.tex
\section{Problem Formulation}

\subsection{Preliminary}

\noindent \textbf{Notation.}
% Throughout this paper, we use the following conventions.
An upper case letter (\textit{e.g.}, $X$) refers to a variable (metric), while a lower case letter (\textit{e.g.}, $x$) represents an \textit{assignment} of the corresponding variable.
By \textit{assignment}, we mean one of the possible values. 
To distinguish variables (values) at different times in a time series, the timestamp will be put on the letter as a superscript.
For example, denote AAS as an upper case letter $Y$, and $y^{(t)}$ refers to the value of AAS at time $t$.
Denote the value range of $Y$ as $Val(Y)$, then we have $y^{(t)} \in Val(Y) = \{0\} \cup \mathbb{R}^{+}$ for the non-negative numeric AAS.
A boldfaced letter means a set of elements (variables or values), \textit{e.g.}, we denote all the metrics as $\mathbf{V}$ while $\mathbf{v}$ is an assignment of $\mathbf{V}$.

% Throughout this work, we use the following conventions.
% An upper case letter (\textit{e.g.}, $X$) refers to a variable (a metric), while a lower case letter (\textit{e.g.}, $x$) represents an assignment of the corresponding variable.
% A bold-face letter means a set of elements, \textit{e.g.}, we denote all the metrics as $\mathbf{V}$ while $\mathbf{v}$ is an assignment of $\mathbf{V}$.
% Possible values of $\mathbf{X}$ are written as $Val(\mathbf{X})$, \textit{i.e.}, $\mathbf{x} \in Val(\mathbf{X})$.
% To distinguish variables (values) in a time-series, the timestamp will be put as a superscript, \textit{e.g.}, $x^{(t)} \in Val^{(t)}(X)$.

\noindent \textbf{The Ladder of Causation.}
% Judea Pearl discovers and studies the ``Ladder of Causation''~\cite{Pearl:2018}.
We formulate the problem with Judea Pearl's ``Ladder of Causation''~\cite{Bareinboim:2022}.
The first layer of the causal ladder encodes the observational knowledge $\mathcal{L}_{1}(\mathbf{V}) = P(\mathbf{V})$, where $P(\mathbf{V})$ is a joint probability distribution.
% \ykl{, where P(V) means ....}
% Meanwhile, the second layer encodes the interventional knowledge $\mathcal{L}_{2}(\mathbf{V}, \mathbf{m}) = P_{\mathbf{m}}(\mathbf{V})$, where $P_{\mathbf{m}}(\mathbf{V}) = P(\mathbf{V} \mid do(\mathbf{m}))$ and $\mathbf{M} \subseteq\mathbf{V}$.
Meanwhile, the second layer encodes the interventional knowledge $\mathcal{L}_{2}(\mathbf{m}) = P_{\mathbf{m}}$, where $P_{\mathbf{m}}(\mathbf{V}) = P(\mathbf{V} \mid do(\mathbf{m}))$ and $\mathbf{M} \subseteq\mathbf{V}$.
% $P(\mathbf{V} \mid do(\mathbf{m}))$ describes the probability distribution of $\mathbf{V}$ when $\mathbf{M}$ is fixed to the given value $\mathbf{m}$, where $do$ is named \textit{do-operator}~\cite{Pearl:2009}, representing an intervention.
The \textit{do-operator} $do(\mathbf{m})$ means fixing variables $\mathbf{M}$ to the given values $\mathbf{m}$, also called an \textit{intervention}~\cite{Pearl:2009}.
So that $P(\mathbf{V} \mid do(\mathbf{m}))$ indicates the probability distribution over $\mathbf{V}$ under the intervention to $\mathbf{M}$.
Finally, the third layer encodes the counterfactual knowledge, reasoning about what if another situation happened in the past.
For example, it requires the counterfactual knowledge to predict the latency with sufficient computing resources when high latency and full CPU usage are observed.
The hierarchy of the causal ladder \textit{almost never} collapses (named CHT, Causal Hierarchy Theorem~\cite{Bareinboim:2022}).
If we want to answer the question at Layer i, we need knowledge at Layer i or higher~\cite{Bareinboim:2022}.
% The ladder of causation defines a hierarchy that \textit{almost never} collapses, named Causal Hierarchy Theorem (CHT)~\cite{Bareinboim:2022}.
% In other words, ``to answer questions at Layer i, one needs knowledge at Layer i or higher.''~\cite{Bareinboim:2022}
% \ykl{The hierarchy of the Ladder of Causation(named CHT, Causal Hierarchy Theorem ~\cite{Bareinboim:2022}) \textit{almost never} collapses. When we want to answer the question at Layer i, we have to needs knowledge at Layer i or higher.~\cite{Bareinboim:2022}}

\noindent \textbf{Structural Causal Model (SCM).}
We model the relations among metrics via the structural causal model~\cite{Pearl:2009}.
An SCM contains a set of structural equations shown in \eqnref{eqn:scm}, where $V_{i} \in \mathbf{V}$ and $\mathbf{Pa}(V_{i}) \subseteq \mathbf{V}$.
% The parameters in \eqnref{eqn:scm} can be divided into two parts: 1) observed variables $\mathbf{Pa}(V_{i})$, named parents (direct causes) of $V_{i}$, and 2) unobserved ones $\mathbf{u}_i$, \textit{e.g.}, noises that we are not interested in.
\eqnref{eqn:scm} contains two kinds of parameters:
1) assignments of observed variables $\mathbf{Pa}(V_{i})$, named parents (direct causes) of $V_{i}$, and
2) assignments of unobserved variables $\mathbf{U}_i$, where $\mathbf{U}_{i} \cap \mathbf{V} = \emptyset$.

\begin{equation}\label{eqn:scm}
	v_{i} = f_{i}(\mathbf{pa}(V_{i}), \mathbf{u}_{i})
\end{equation}

Denote the graph encoded by the SCM as $\mathcal{G} = (\mathbf{V}, \mathbf{E})$, where $\mathbf{E} = \{V_{j} \to V_{i} \mid V_{j} \in \mathbf{Pa}(V_{i})\}$ is the set of directed edges.
In contrast to $\mathbf{Pa}$, $\mathbf{Ch}(V_{i}) = \{V_{j} \mid V_{i} \in \mathbf{Pa}(V_{j})\}$ represents the children of $V_{i}$.
This work rests on the following assumptions.
\begin{description}
    \item[DAG] $\mathcal{G}$ is a directed acyclic graph (DAG)~\cite{Pearl:2009}, following related works in OSS~\cite{Chen:2014,Wang:2018,Gan:2021}.
    \item[Markovian] ``The exogenous parent sets $\mathbf{U}_{i}, \mathbf{U}_{j}$ are independent whenever $i \neq j$''~\cite{Bareinboim:2022}, \textit{i.e.}, $(\forall i \neq j) \mathbf{U}_{i} \ci \mathbf{U}_{j}$, where $\ci$ means independent.
    % \item[Causal Sufficiency] ``All common drivers are among the observed variables''~\cite{Runge:2019}, \textit{i.e.}, $(\forall i \neq j) \mathbf{U}_{i} \ci \mathbf{U}_{j}$, where $\ci$ means independent.
    % \item[Causal Sufficiency] The SCM under discussion is \textit{Markovian}~\cite{Bareinboim:2022} (also known as the \textit{Causal Sufficiency} assumption~\cite{Runge:2019}), \textit{i.e.}, $(\forall i \neq j) \mathbf{U}_{i} \ci \mathbf{U}_{j}$, where $\ci$ means independent;
    \item[Faithfulness~\cite{Pearl:2009}] Any intervention makes an observable change, \textit{i.e.}, $P(V_{i} \mid \mathbf{pa}(V_{i}), do(v_{i})) \neq P(V_{i} \mid \mathbf{pa}(V_{i}))$.
\end{description}
Under the DAG assumption and the Markovian assumption, $\mathcal{G}$ can be taken as a CBN~\cite{Bareinboim:2022}.

% Denote the graph encoded by the SCM as $\mathcal{G} = (\mathbf{V}, \mathbf{E})$, where $\mathbf{E} = \{V_{j} \to V_{i} \mid V_{j} \in \mathbf{Pa}(V_{i})\}$ is the set of directed edges.
% % \ykl{With SCM, realations between two variables can be encoded as a directed edge $\mathbf{E} = \{V_{j} \to V_{i} \mid V_{j} \in \mathbf{Pa}(V_{i})\}$, and finally a causation graph $\mathcal{G} = (\mathbf{V}, \mathbf{E})$ is formed.}
% In contrast to $\mathbf{Pa}$, $\mathbf{Ch}(V_{i}) = \{V_{j} \mid V_{i} \in \mathbf{Pa}(V_{j})\}$ represents the children of $V_{i}$.
% We assume that $\mathcal{G}$ is a directed acyclic graph (DAG)~\cite{Pearl:2009}, following \cite{Chen:2014,Wang:2018,Gan:2021} in OSS.
% % We adopt the common assumption that $\mathcal{G}$ is a directed acyclic graph (DAG) for a real application.
% In the DAG, denote ancestors of $V_{i}$ as $\mathbf{An}(V_{i}) = \{V_{i}\} \cup \bigcup_{X \in \mathbf{An}(V_{i})} \mathbf{Pa}(X)$.
% % \ykl{And we denote ... in DAG}
% %In contrast, $\mathbf{Desc}(V_{i}) = \{V_{j} | \left(\exists X \in \{V_{i}\} \cup \mathbf{Desc}(V_{i})\right) V_{j} \in \mathbf{Pa}(X) \}$ represents the descendants of a node $V_{i}$.
% This work further rests on the following standard assumptions~\cite{Runge:2019}.
% Readers can find a thorough discussion about these assumptions in \cite{Pearl:2009}.
% Under the following assumptions, $\mathcal{G}$ can be taken as a CBN~\cite{Bareinboim:2022}.

% \begin{itemize}
% 	\item \textbf{Causal Sufficiency} implies that ``all common drivers are among the observed variables''~\cite{Runge:2019}, \textit{i.e.}, $(\forall i \neq j) \mathbf{U}_{i} \ci \mathbf{U}_{j}$, where $\ci$ means independent.
% % 	\ykl{What is the meaning of this example?}\mingjie{The content following ``\textit{i.e.}'' is the mathematical expression of Causal Sufficiency, not example}
% 	\item \textbf{Causal Markov Condition} implies that $V_{i} \ci V_{j} \mid \mathbf{Pa}(V_{i})$, where $V_{i} \in \mathbf{V}$ and $V_{j} \in \mathbf{An}(V_{i}) \setminus \{V_{i}\} \setminus \mathbf{Pa}(V_{i})$.
% 	\item \textbf{Faithfulness} requires that any intervention makes an observable change, \textit{i.e.}, $P(V_{i} \mid \mathbf{pa}(V_{i}), do(v_{i})) \neq P(V_{i} \mid \mathbf{pa}(V_{i}))$.
% \end{itemize}

\subsection{Root Cause Analysis and Causal Inference}

We set up a concept mapping between the RCA problem and causal inference.
\begin{itemize}
    \item A fault in OSS is mapped to an unexpected intervention;
    \item Fault-free data come from the observational distribution;
    \item Faulty data come from an interventional distribution.
\end{itemize}
Based on the mapping above, we define a new causal inference task as \textit{intervention recognition} (Definition~\ref{def:intervention-recognition}).
We formulate RCA discussed in this work as an intervention recognition task in OSS.

\begin{definition}[Intervention Recognition, IR]\label{def:intervention-recognition}
    For a given SCM $\mathcal{M}$, let $\mathcal{L}_{1}$ be the observational distribution of $\mathcal{M}$ and $P_{m} = P(\mathbf{V} \mid do(\mathbf{m}))$ be the interventional distribution of a certain intervention $do(\mathbf{m})$.
    \textit{Intervention recognition} is to find $\mathbf{m}$ based on $\mathcal{L}_{1}$ and $P_{m}$.
\end{definition}

\begin{definition}[Root Cause]\label{def:root-cause}
    The \textit{root cause} is the intervened variables ($\mathbf{M}$).
    Each element of $\mathbf{M}$ is named a \textit{root cause variable}.\footnote{We also use \textit{root cause indicator} and \textit{root cause metric} according to the context.}
\end{definition}

\section{Intervention Recognition Criterion}\label{sec:rc-criterion}

% 明杰：
%   我觉得“argue”比“prove”更合适
%   Theorem 3.1的成立是有条件的，正文中的这句话没有交代条件
%   把条件写在正文里会更难懂
We argue that IR shall be positioned at the second layer in the ladder of causation, as shown in \thmref{thm:rca-l2}.
The proof of \thmref{thm:rca-l2} is provided in Appendix~\ref{sec:proof-rca-l2}.
The key to the proof is that IR is the inverse mapping of $\mathcal{L}_{2}$ under the adopted assumptions.
Combining \thmref{thm:rca-l2} with CHT~\cite{Bareinboim:2022}, we further obtain Corollary~\ref{thm:l2-necessary} and \ref{thm:l3-unnecessary}.

\begin{theorem}\label{thm:rca-l2}
	For a given SCM $\mathcal{M}$ with a CBN $\mathcal{G}$, the knowledge of IR for $\mathcal{M}$ is equivalent to $\mathcal{L}_{2}$ under the Faithfulness assumption.
\end{theorem}

\begin{corollary}\label{thm:l2-necessary}
	We need the knowledge at Layer 2 (interventional) to conduct IR.
\end{corollary}

\begin{corollary}\label{thm:l3-unnecessary}
	The knowledge at Layer 3 (counterfactual) is not necessary to conduct IR.
\end{corollary}

Hence, we propose to take full advantage of the CBN, as the CBN is a known bridge between observational data and interventional knowledge~\cite{Bareinboim:2022}.
We argue that \thmref{thm:rc-criterion} is a necessary and sufficient condition for a variable to be intervened.
The proof of \thmref{thm:rc-criterion} is provided in Appendix~\ref{sec:proof-rc-criterion}.
Based on our concept mapping between RCA and causal inference, the \nameref{thm:rc-criterion} is also a criterion to find root cause indicators.

\begin{theorem}[Intervention Recognition Criterion]\label{thm:rc-criterion}
	Let $\mathcal{G}$ be a CBN and $\mathbf{Pa}(V_{i})$ be the parents of $V_{i}$ in $\mathcal{G}$.
	Under the Faithfulness assumption, $V_{i}$ is intervened \textit{iff} $V_{i}$ no longer follows the distribution defined by $\mathbf{pa}(V_{i})$, \textit{i.e.},
	\begin{equation*}
		V_{i} \in \mathbf{M} \Leftrightarrow P_{\mathbf{m}}(V_{i} \mid \mathbf{pa}(V_{i})) \neq \mathcal{L}_{1}(V_{i} \mid \mathbf{pa}(V_{i}))
	\end{equation*}
\end{theorem}

% When a fault occurs and affects $V_{i}$ directly (\textit{i.e.}, $V_{i} \in \mathbf{M}$), $v_{i}$ is no longer defined by $f_{i}$ in \eqnref{eqn:scm}\footnote{We refer readers to the Manipulation Theorem~\cite{Spirtes:1993} for more discussion.}.
% There are only two basic situations:
% 1) $f_{i}$ is replaced by $f_{i}^{\prime}$, or
% 2) the distribution of unobserved variables changed from $P(\mathbf{U}_{i})$ to $P^{\prime}(\mathbf{U}_{i})$.
% It is a special case of the first situation that $\mathbf{Pa}(V_{i})$ (or $\mathbf{U}_{i}$) is changed to $\mathbf{Pa}^{\prime}(V_{i})$ (or $\mathbf{U}_{i}^{\prime}$).
% If only $\mathbf{pa}(V_{i})$ (the value of $\mathbf{Pa}(V_{i})$ ) changes, $V_{i}$ should not be considered for the root cause, as the failure just propagates through $V_{i}$.
% In contrast, $V_{i}$ is the only visible indicator of $\mathbf{U}_{i}$ under the assumption of Causal Sufficiency.
% Hence, $V_{i}$ is the desired result if $\mathbf{u}_{i}$ changes.
% In summary, we argue that the following criterion (\thmref{thm:rc-criterion}) is a necessary and sufficient condition for a metric to be an indicator of the root cause.
% The proof of \thmref{thm:rc-criterion} is provided in Appendix~\ref{sec:proof-rc-criterion}.
